% Chapter 9: 估值模型
\section{板块估值锚（市场共识PE，2026年6月）}

\begin{table}[H]
\centering\small
\renewcommand{\arraystretch}{1.2}
\begin{tabularx}{\textwidth}{l c X}
\toprule
\textbf{板块} & \textbf{共识PE} & \textbf{参考依据} \\
\midrule
人形机器人核心零部件 & 50--60$\times$ & Wind人形机器人指数PE 56.89$\times$ \\
AI算力/服务器 & 30--40$\times$ & 协创数据类NVIDIA供应链高增长 \\
智能驾驶Tier1 & 30--40$\times$ & 德赛西威历史中枢30--35$\times$，Thor预期40$\times$ \\
AI边缘计算 & 35--50$\times$ & Jetson生态高增长 \\
传感器/执行器 & 40--55$\times$ & 六维力/空心杯稀缺赛道溢价 \\
\bottomrule
\end{tabularx}
\caption{板块估值锚}
\end{table}

\section{逐标的估值测算}

\begin{table}[H]
\centering
\scriptsize
\renewcommand{\arraystretch}{1.2}
\begin{tabularx}{\textwidth}{l c c c c c c c}
\toprule
\textbf{标的} & \textbf{现价/市值} & \textbf{Q1净利} & \textbf{2026E全年} & \textbf{PE} & \textbf{目标市值} & \textbf{目标价} & \textbf{空间} \\
\midrule
德赛西威 & 90.8/541亿 & 4.61亿(归母) & 25--28亿 & 35$\times$ & 875--980亿 & 147--164 & \textcolor{riskgreen}{\textbf{+62\textasciitilde81\%}} \\
协创数据 & 213/1035亿 & 7.5亿 & 25--30亿 & 35$\times$ & 875--1050亿 & 180--216 & -5\textasciitilde+2\% \\
天准科技 & 101/196亿 & 0.047亿 & PS 10$\times$ & — & 200亿 & 103 & +2\% \\
柯力传感 & 70/197亿 & 0.41亿 & 3.5--4.5亿 & 50$\times$ & 175--225亿 & 62--80 & -11\textasciitilde+14\% \\
绿的谐波 & 420/770亿 & 0.87亿 & 1.8--2.0亿 & 55$\times$ & 100--112亿 & 55--61 & \textcolor{riskred}{\textbf{-87\textasciitilde-85\%}} \\
鸣志电器 & 66/277亿 & 0.14亿 & 1--1.5亿 & 50$\times$ & 50--75亿 & 12--18 & \textcolor{riskred}{\textbf{-73\textasciitilde-82\%}} \\
奥比中光 & 148/280亿 & 0.31亿 & 1.5--2亿 & 55$\times$ & 82--110亿 & 44--59 & \textcolor{riskred}{\textbf{-60\textasciitilde-70\%}} \\
索辰科技 & 236/200亿 & 亏损 & PS 20$\times$ & — & 100亿 & 118 & \textcolor{riskred}{\textbf{-50\%}} \\
中大力德 & 89/174亿 & 亏损 & PS 8$\times$ & — & 96亿 & 49 & \textcolor{riskred}{\textbf{-45\%}} \\
\bottomrule
\end{tabularx}
\caption{逐标的估值测算（2026E利润$\times$共识PE）}
\end{table}

\section{三档估值体系}

\begin{keyinsight}[三档估值体系说明]
\begin{itemize}[leftmargin=1.5em]
  \item \textcolor{riskamber}{\textbf{乐观档}} = 券商远期折现逻辑（行业按最乐观剧本爆发）
  \item \textcolor{blue}{\textbf{中性档}} = 2026E业绩兑现 $\times$ 合理PE（安全边际锚）
  \item \textcolor{riskgreen}{\textbf{悲观档}} = 叙事破裂+不及预期的防御线
\end{itemize}
\end{keyinsight}

\begin{table}[H]
\centering
\scriptsize
\renewcommand{\arraystretch}{1.15}
\begin{tabularx}{\textwidth}{l c c c c c X}
\toprule
\textbf{标的} & \textbf{现价} & \textbf{\textcolor{riskgreen}{悲观}} & \textbf{\textcolor{blue}{中性}} & \textbf{\textcolor{riskamber}{乐观}} & \textbf{泡沫度} & \textbf{策略} \\
\midrule
德赛西威 & 90.8 & 105 & 147--164 & 180--200 & -41\%低估 & \textbf{首选加仓} \\
协创数据 & 213 & 129 & 180--216 & 277 & -1\textasciitilde+18\% & 基本到位，超预期加仓 \\
天准科技 & 101 & 62 & 103--124 & 155 & -19\textasciitilde-2\% & 接近中性，等放量 \\
智微智能 & 100 & 44 & 78--100 & 125 & +0\textasciitilde+28\% & 偏高，需Q2超预期 \\
柯力传感 & 70 & 31 & 62--80 & 96 & -12\textasciitilde+13\% & 中性，等7月Optimus \\
绿的谐波 & 420 & 27 & 55--61 & 326/800 & +589\textasciitilde664\% & \textcolor{riskred}{极端泡沫，止盈} \\
鸣志电器 & 66 & 9 & 12--18 & 50 & +267\textasciitilde450\% & \textcolor{riskred}{严重泡沫} \\
奥比中光 & 148 & 26 & 44--59 & 180 & +151\textasciitilde236\% & 泡沫区，关注订单 \\
索辰科技 & 236 & 59 & 118 & 472 & +100\% & 纯概念，需产品验证 \\
中大力德 & 89 & 18 & 49 & 200 & +82\% & 赌宇树IPO，风险高 \\
\bottomrule
\end{tabularx}
\caption{三档估值体系（全标的）}
\end{table}

\section{估值倒挂深度解析}

以绿的谐波为例：我们按「2026E净利$\times$55x行业PE」算出合理价55--61元，但券商一致目标价中枢318--326元（仍低于现价420元）。三种完全不同的估值框架在对话：

\begin{table}[H]
\centering\scriptsize
\renewcommand{\arraystretch}{1.15}
\begin{tabularx}{\textwidth}{l X l l c}
\toprule
\textbf{框架} & \textbf{算法} & \textbf{使用者} & \textbf{隐含假设} & \textbf{绿的目标价} \\
\midrule
\textcolor{blue}{本报告（中性）} & 2026E静态利润$\times$行业PE & 买方风控 & 今年业绩兑现 & 55--61元 \\
\textcolor{riskamber}{券商（乐观）} & 2028--30E远期利润$\times$成熟PE→折现 & 卖方研究员 & 3年后按剧本爆发 & 318--326元 \\
\textcolor{riskred}{市场（情绪）} & 2026E利润$\times$100--400x & 游资/散户 & 叙事热度持续 & 420元 \\
\bottomrule
\end{tabularx}
\caption{三种估值框架对比}
\end{table}

\textbf{券商目标价326元反推}：2028E Optimus 50万台 $\times$ 14个减速器 $\times$ ASP 1500元 $\times$ 70\%市占 = 收入73.5亿 $\times$ 25\%净利率 = 18.4亿 $\times$ 25x = 460亿 $\div$ 折现1.25x = 368亿 $\div$ 1.83亿股 $\approx$ 200元/股（给溢价后到326元）。\textit{本质等价：2亿(2026E)$\times$184xPE}。

\begin{riskbox}[卖方目标价系统性上偏]
券商给了乐观目标价3年后没达到无人追责；给了保守目标价导致客户踏空则被骂。卖方目标价天然有\textbf{上偏50\%\textasciitilde100\%}的统计倾向。本报告估值是买方风控视角：告诉你"叙事破裂后股价会回落到哪里"——安全边际，不是天花板。
\end{riskbox}

\section{相对估值横向对比（PEG七维表）}

\begin{table}[H]
\centering
\scriptsize
\renewcommand{\arraystretch}{1.1}
\begin{tabularx}{\textwidth}{l c c c c c c c c}
\toprule
\textbf{标的} & \textbf{市值} & \textbf{PE} & \textbf{PS} & \textbf{PB} & \textbf{增速} & \textbf{PEG} & \textbf{机构分} & \textbf{结论} \\
\midrule
德赛西威 & 541亿 & 21$\times$ & 1.3$\times$ & 4.2 & +5--15\% & 1.5--4.2 & A(94) & \textcolor{blue}{合理偏低} \\
协创数据 & 1035亿 & 38$\times$ & 6.9$\times$ & 12 & +180\% & \textbf{0.17--0.21} & B+(72) & \textcolor{riskgreen}{PEG最低} \\
工业富联 & 8300亿 & 19$\times$ & 1.0$\times$ & — & +50\% & \textbf{0.32--0.48} & A+(100) & \textcolor{riskgreen}{防御首选} \\
中际旭创 & 1600亿 & 29$\times$ & 3.4$\times$ & — & +40--70\% & 0.41--0.73 & A+(98) & \textcolor{riskgreen}{低估} \\
绿的谐波 & 770亿 & 397$\times$ & 154$\times$ & 26 & +50--70\% & \textbf{5.7--7.9} & D(41) & \textcolor{riskred}{全表最贵} \\
鸣志电器 & 277亿 & 222$\times$ & 24.5$\times$ & 6.2 & -20--+25\% & $>$8.9 & D(35) & \textcolor{riskred}{极端泡沫} \\
奥比中光 & 280亿 & 160$\times$ & 15.6$\times$ & 6.4 & +200--300\% & 0.53--0.80 & D+(47) & \textcolor{riskamber}{PEG低但机构不买} \\
索辰科技 & 200亿 & 亏损 & 66.7$\times$ & 11 & 营收+50--70\% & PSG 0.29 & D-(15) & \textcolor{riskred}{PS全表第二高} \\
天准科技 & 196亿 & 扭亏 & 16.3$\times$ & 5.8 & 营收+60--80\% & PSG 0.15 & D+(45) & \textcolor{blue}{PSG低估待解禁} \\
中大力德 & 174亿 & 亏损 & 14.5$\times$ & 6.9 & 营收+33\% & PSG 0.24 & D(32) & \textcolor{riskamber}{基本面烂纯赌催化} \\
\bottomrule
\end{tabularx}
\caption{十大核心标的PEG/PS/机构可投性七维横向对比}
\end{table}

\textbf{三大洞察}：① PEG最便宜前3：协创(0.19)/工业富联(0.40)/中际旭创(0.57)——大资金底仓；② PEG最贵前3：鸣志($>$8.9)/绿的谐波(6.8)——游资炒作；③ 三维共振（PEG低+机构友好+有空间）前5 = 机构核心池。

\section{"炒预期"的正确姿势}

\begin{table}[H]
\centering\scriptsize
\renewcommand{\arraystretch}{1.15}
\begin{tabularx}{\textwidth}{l c c c X}
\toprule
\textbf{标的} & \textbf{情绪PE上限} & \textbf{情绪目标价} & \textbf{vs现价} & \textbf{什么条件能到} \\
\midrule
德赛西威 & 45$\times$ & 189元 & +108\% & Q2环比回升+Thor首批定点公告 \\
协创数据 & 45$\times$ & 277元 & +30\% & Q2净利$\geq$8亿+Coalition大单 \\
天准科技 & PS 15$\times$ & 155元 & +53\% & Q2具身智能域控批量出货 \\
绿的谐波 & 250$\times$ & 278元 & \textcolor{riskred}{-34\%} & 即使250xPE目标也低于现价！ \\
柯力传感 & 60$\times$ & 96元 & +37\% & 7月Optimus投产+月出货3000只 \\
\bottomrule
\end{tabularx}
\caption{"炒预期"情绪PE上限测算}
\end{table}

\section{Q2--Q3核心催化事件时间线}

\begin{figure}[H]
\centering
\begin{tikzpicture}[
    x=2.2cm,
    event/.style={font=\scriptsize, rounded corners=2pt, minimum height=0.55cm, inner sep=3pt, anchor=west},
    track label/.style={font=\scriptsize\bfseries, anchor=east},
]
% 时间轴
\draw[steelgrey!60, ->, thick] (-0.3, 0) -- (6.8, 0);
\foreach \x/\m in {0/6月, 1/7月, 2/8月, 3/9月, 4/10月, 5/11月, 6/12月} {
    \draw[lightgrey] (\x, -0.1) -- (\x, 4.6);
    \node[below, font=\tiny, steelgrey] at (\x, 0) {\m};
}

% Track 1: Optimus
\node[track label] at (-0.4, 4) {Optimus};
\node[event, fill=nvgreen!15, draw=nvgreen!50] at (0, 4) {\tiny 产线改造};
\node[event, fill=nvgreen!25, draw=nvgreen!60] at (1.5, 4) {\tiny V3正式投产};
\node[event, fill=nvgreen!15, draw=nvgreen!50] at (3, 4) {\tiny 首批量产交付};

% Track 2: Cosmos 3
\node[track label] at (-0.4, 3) {Cosmos 3};
\node[event, fill=navy!15, draw=navy!40] at (0, 3) {\tiny GTC发布\checkmark};
\node[event, fill=navy!25, draw=navy!50] at (1, 3) {\tiny Coalition落地};
\node[event, fill=navy!15, draw=navy!40] at (3.5, 3) {\tiny 商业化订单};

% Track 3: 宇树
\node[track label] at (-0.4, 2) {宇树IPO};
\node[event, fill=riskamber!15, draw=riskamber!50] at (1, 2) {\tiny 科创板上市};
\node[event, fill=riskamber!25, draw=riskamber!60] at (3, 2) {\tiny H2+规模化量产};

% Track 4: 财报
\node[track label] at (-0.4, 1) {二季报};
\node[event, fill=steelgrey!15, draw=steelgrey!40] at (1.5, 1) {\tiny 半年报预约};
\node[event, fill=steelgrey!15, draw=steelgrey!40] at (3, 1) {\tiny 密集披露期};

% Q4 DRIVE Thor
\node[event, fill=navy!30, draw=navy!70, font=\scriptsize\bfseries] at (4.5, 3.5) {\tiny DRIVE Thor量产};
\end{tikzpicture}
\caption{2026年Q2--Q4核心催化事件时间线}
\end{figure}

\section{核心标的Q2业绩观察门槛}

以下为各标的Q2需达到的\textbf{最低业绩门槛}——低于此则当前估值不可支撑，高于此则有上修空间：

\begin{table}[H]
\centering\scriptsize
\renewcommand{\arraystretch}{1.15}
\begin{tabularx}{\textwidth}{l c c X X}
\toprule
\textbf{标的} & \textbf{Q1实际} & \textbf{Q2门槛} & \textbf{门槛逻辑} & \textbf{若超预期} \\
\midrule
德赛西威 & 4.61亿(归母) & $\geq$6.5亿 & Q1为全年低点（价格战+春节），Q2环比回升确认复苏斜率 & $\geq$8亿→年化28亿→+81\%空间确认 \\
协创数据 & 7.5亿 & $\geq$8亿 & 一致预期25--30亿需Q2--Q4均$\geq$6亿；超8亿→年化32亿 & $\geq$9亿→年化36亿→空间+30\% \\
智微智能 & 1.09亿 & $\geq$1.2亿 & 当前100元隐含年化5亿$\times$50x；Q2需环比持续1.2亿+ & $<$1亿→当前价-22\%回调压力 \\
柯力传感 & 0.41亿(含公允拖累) & 扣非$\geq$0.8亿 & 7月Optimus投产带量，扣非环比翻倍是验证信号 & 扣非$\geq$1.2亿→年化4.5亿→80元(+14\%) \\
天准科技 & 0.047亿(扭亏) & $\geq$0.5亿 & 具身域控Q2放量确认，亏损不可再现 & $\geq$1亿→全年5亿可期→PS 12x=240亿 \\
奥比中光 & 0.31亿 & $\geq$0.5亿 & 机器人3倍增速持续性验证 & $\geq$0.8亿→全年2.5亿→PE 112x仍高 \\
\bottomrule
\end{tabularx}
\caption{Q2业绩观察门槛（低于此=估值不可支撑）}
\end{table}

\section{德赛西威DRIVE Thor业绩兑现路线图}

\begin{table}[H]
\centering\small
\renewcommand{\arraystretch}{1.2}
\begin{tabularx}{\textwidth}{l l X c}
\toprule
\textbf{时间节点} & \textbf{里程碑} & \textbf{业绩贡献} & \textbf{确定性} \\
\midrule
2025.3 & IPU14Pro全球首个SOP上车（iCAR 03T） & 验证性订单，收入$<$5亿 & \textcolor{riskgreen}{已完成} \\
2026Q1 & Thor域控正式量产，10+车企定点确认 & 营收贡献约15--20亿（含NRE） & \textcolor{riskgreen}{已完成} \\
2026Q2--Q3 & 极氪/智己/理想等头部车企批量交付 & 环比放量，预计Q2单季贡献35--45亿 & \textcolor{riskamber}{进行中} \\
2026Q4 & 首批累计上车10万辆（官方指引） & 10万辆$\times$ASP 3万元=\textbf{30亿直接域控收入} & \textcolor{riskamber}{指引} \\
2027全年 & Thor全车型全年30--50万辆规模化 & \textbf{全年Thor相关营收合计约240亿}（含域控硬件+软件授权+项目交付费） & 远期 \\
\midrule
\multicolumn{4}{l}{\textit{240亿拆解：30--50万辆$\times$ASP 3万/车=90--150亿（硬件）+ DRIVE软件年度授权50亿 + 项目服务费40亿}} \\
\bottomrule
\end{tabularx}
\caption{德赛西威DRIVE Thor业绩兑现时间表}
\end{table}

\begin{keyinsight}[德赛西威核心判断]
Q1净利-21\%是\textbf{布局年阵痛}（Thor研发投入高峰+价格战冲击），非趋势恶化。H2起随Thor量产爬坡，业绩弹性为全表最大（单里程碑+240亿=协创3个里程碑之和）。\textbf{2026Q4累计10万辆}是最核心验证节点——达成则+107--142\%空间完全打开，不达成（$<$3万辆）则需下修30\%。
\end{keyinsight}

\section{估值结论}

\begin{table}[H]
\centering\small
\renewcommand{\arraystretch}{1.2}
\begin{tabularx}{\textwidth}{c l c X}
\toprule
\textbf{分级} & \textbf{标的} & \textbf{空间} & \textbf{核心逻辑} \\
\midrule
最优 & 德赛西威 & +62\textasciitilde81\% & Q1全年低点，Thor量产+130亿订单释放 \\
等催化 & 天准/柯力 & 0\textasciitilde+23\% & 具身智能放量/Optimus投产 \\
到位 & 协创数据 & -5\textasciitilde+8\% & 一致预期基本price-in \\
\textcolor{riskred}{透支} & 绿的谐波 & -85\% & 770亿vs合理112亿=7倍透支 \\
\textcolor{riskred}{透支} & 鸣志/奥比/索辰/中大 & -45\textasciitilde-82\% & 纯情绪驱动 \\
\bottomrule
\end{tabularx}
\caption{估值分级结论}
\end{table}
